\begin{abstract}
Reforms designed to strengthen electoral integrity can also make electoral participation more administratively demanding. This paper studies that tradeoff through the staggered rollout of biometric voter registration across Brazilian municipalities between 2008 and 2018. Using administrative electoral records, I show that biometric registration reduced the size of the registered electorate at adoption and disproportionately lowered the share of less-educated voters on the rolls. I then combine those administrative results with LAPOP survey data to examine whether exposure to the reform is associated with changes in trust in institutions, democratic attitudes, and perceptions of electoral integrity. The survey evidence suggests higher trust in institutions among lower-education respondents in treated municipalities, while effects on broader democratic attitudes are weaker and evidence of backlash among Bolsonaro-aligned respondents is limited. The results highlight a central tension of electoral modernization: reforms intended to secure the voter registry can also reshape who remains represented within it.
\end{abstract}
